\chapter{1991年三南卷}
\section{单选题}

\begin{problem}
\(\sin 15^{\circ}\cos 30^{\circ}\sin 75^{\circ}\) 的值等于（\quad）

\choices
    {\(\frac{\sqrt{3}}{4}\)}
    {\(\frac{\sqrt{3}}{8}\)}
    {\(\frac{1}{8}\)}
    {\(\frac{1}{4}\)}

\end{problem}

\begin{answer}
B
\end{answer}

\begin{solution}
由积化和差公式，\(\sin 15^{\circ}\sin 75^{\circ}=\frac{\cos 60^{\circ}-\cos 90^{\circ}}{2}=\frac{1}{4}\)，所以 \(\sin 15^{\circ}\cos 30^{\circ}\sin 75^{\circ}=\frac{1}{4}\cdot\frac{\sqrt{3}}{2}=\frac{\sqrt{3}}{8}\)，故选 B.
\end{solution}

\begin{problem}
已知一个等差数列的第 \(5\) 项等于 \(10\)，前 \(3\) 项的和等于 \(3\). 那么（\quad）

\choices
    {它的首项是 \(-2\)，公差是 \(3\)}
    {它的首项是 \(2\)，公差是 \(-3\)}
    {它的首项是 \(-3\)，公差是 \(2\)}
    {它的首项是 \(3\)，公差是 \(-2\)}

\end{problem}

\begin{answer}
A
\end{answer}

\begin{solution}
设等差数列的首项为 \(a_1\)，公差为 \(d\). 由题意，\(a_1+4d=10\)，且 \(S_3=3(a_1+d)=3\)，所以 \(a_1+d=1\). 两式相减得 \(3d=9\)，从而 \(d=3\)，\(a_1=-2\)，故选 A.
\end{solution}

\begin{problem}
设正六棱锥的底面边长为 \(1\)，侧棱长为 \(\sqrt{5}\)，那么它的体积为（\quad）

\choices
    {\(6\sqrt{3}\)}
    {\(2\sqrt{3}\)}
    {\(\sqrt{3}\)}
    {\(2\)}

\end{problem}

\begin{answer}
C
\end{answer}

\begin{solution}
正六边形的外接圆半径等于其边长，因此底面中心到任一顶点的距离为 \(1\). 设正六棱锥的高为 \(h\)，由侧棱长为 \(\sqrt{5}\) 得 \(h^2+1^2=5\)，所以 \(h=2\). 底面面积为 \(6\cdot\frac{\sqrt{3}}{4}=\frac{3\sqrt{3}}{2}\)，故体积为 \(\frac{1}{3}\cdot\frac{3\sqrt{3}}{2}\cdot 2=\sqrt{3}\)，选 C.
\end{solution}

\begin{problem}
在直角坐标系 \(xOy\) 中，参数方程 \(\begin{cases}
x=2t+1,\\
y=2t^2-1,
\end{cases}\)（其中 \(t\) 是参数）表示的曲线是（\quad）

\choices
    {双曲线}
    {抛物线}
    {直线}
    {圆}

\end{problem}

\begin{answer}
B
\end{answer}

\begin{solution}
由参数方程 \(x=2t+1\) 得 \(t=\frac{x-1}{2}\)，代入第二式得 \(y=2\left(\frac{x-1}{2}\right)^2-1=\frac{1}{2}(x-1)^2-1\). 这是顶点为 \((1,-1)\) 的抛物线，故选 B.
\end{solution}

\begin{problem}
设全集为自然数集 \(N\)，\(E=\{x\mid x=2n, n\in \N\}\)，\(F=\{x\mid x=4n, n\in \N\}\)，那么集合 \(N\) 可以表示成（\quad）

\choices
    {\(E\cap F\)}
    {\(\overline{E}\cup F\)}
    {\(E\cup \overline{F}\)}
    {\(\overline{E}\cap \overline{F}\)}

\end{problem}

\begin{answer}
C
\end{answer}

\begin{solution}
因为每个 \(4\) 的倍数也是 \(2\) 的倍数，所以 \(F\subset E\). 对任意 \(x\in N\)，若 \(x\in F\)，则 \(x\in E\)；若 \(x\notin F\)，则 \(x\in\overline{F}\). 因而 \(N=E\cup\overline{F}\)，故选 C.
\end{solution}

\begin{problem}
已知 \(z_1\)，\(z_2\) 是两个给定的复数，且 \(z_1\neq z_2\). 它们在复平面上分别对应于点 \(Z_1\) 和点 \(Z_2\). 如果 \(z\) 满足方程 \(\left| z-z_1 \right|-\left| z-z_2 \right|=0\)，那么 \(z\) 对应的点 \(Z\) 的集合是（\quad）

\choices
    {双曲线}
    {线段 \(Z_1Z_2\) 的垂直平分线}
    {分别过 \(Z_1\)，\(Z_2\) 的两条相交直线}
    {椭圆}

\end{problem}

\begin{answer}
B
\end{answer}

\begin{solution}
方程 \(\left|z-z_1\right|-\left|z-z_2\right|=0\) 等价于 \(\left|z-z_1\right|=\left|z-z_2\right|\). 这表示点 \(Z\) 到定点 \(Z_1\)，\(Z_2\) 的距离相等，因此 \(Z\) 在以线段 \(Z_1Z_2\) 为垂直平分线的直线上，故选 B.
\end{solution}

\begin{problem}
设 \(5\pi<\theta<6\pi\)，\(\cos \frac{\theta}{2}=a\)，那么 \(\sin \frac{\theta}{4}\) 等于（\quad）

\choices
    {\(-\frac{\sqrt{1+a}}{2}\)}
    {\(-\frac{\sqrt{1-a}}{2}\)}
    {\(-\sqrt{\frac{1+a}{2}}\)}
    {\(-\sqrt{\frac{1-a}{2}}\)}

\end{problem}

\begin{answer}
D
\end{answer}

\begin{solution}
由 \(5\pi<\theta<6\pi\) 可得 \(\frac{5\pi}{4}<\frac{\theta}{4}<\frac{3\pi}{2}\)，所以 \(\sin\frac{\theta}{4}<0\). 又 \(\sin^2\frac{\theta}{4}=\frac{1-\cos\frac{\theta}{2}}{2}=\frac{1-a}{2}\)，故 \(\sin\frac{\theta}{4}=-\sqrt{\frac{1-a}{2}}\)，选 D.
\end{solution}

\begin{problem}
函数 \(y=\sin x\)，\(x\in \left[\frac{\pi}{2},\frac{3\pi}{2}\right]\) 的反函数为（\quad）

\choices
    {\(y=\arcsin x\)，\(x\in [-1,1]\)}
    {\(y=-\arcsin x\)，\(x\in [-1,1]\)}
    {\(y=\pi+\arcsin x\)，\(x\in [-1,1]\)}
    {\(y=\pi-\arcsin x\)，\(x\in [-1,1]\)}

\end{problem}

\begin{answer}
D
\end{answer}

\begin{solution}
在区间 \(\left[\frac{\pi}{2},\frac{3\pi}{2}\right]\) 上，函数 \(y=\sin x\) 单调递减. 由 \(y=\sin x\) 得 \(x=\pi-\arcsin y\)，交换 \(x\)，\(y\) 得反函数 \(y=\pi-\arcsin x\)，其定义域为原函数的值域 \([-1,1]\)，故选 D.
\end{solution}

\begin{problem}
复数 \(z=-3\left(\sin \frac{4}{3}\pi-\i\cos \frac{4}{3}\pi\right)\) 的辐角的主值是（\quad）

\choices
    {\(\frac{4}{3}\pi\)}
    {\(\frac{5}{3}\pi\)}
    {\(\frac{11}{6}\pi\)}
    {\(\frac{\pi}{6}\)}

\end{problem}

\begin{answer}
C
\end{answer}

\begin{solution}
因为 \(\sin\frac{4\pi}{3}=-\frac{\sqrt{3}}{2}\)，\(\cos\frac{4\pi}{3}=-\frac{1}{2}\)，所以 \(z=-3\left(-\frac{\sqrt{3}}{2}+\frac{\i}{2}\right)=\frac{3\sqrt{3}}{2}-\frac{3}{2}\i\). 复数 \(z\) 位于第四象限，参考角为 \(\frac{\pi}{6}\)，按本题主值取法，其辐角主值为 \(2\pi-\frac{\pi}{6}=\frac{11\pi}{6}\)，故选 C.
\end{solution}

\begin{problem}
满足 \(\sin \left(x-\frac{\pi}{4}\right)\geqslant \frac{1}{2}\) 的 \(x\) 的集合是（\quad）

\choices
    {\(\left\{x\mid 2k\pi+\frac{5\pi}{12}\leqslant x\leqslant 2k\pi+\frac{13\pi}{12}, k\in \Z\right\}\)}
    {\(\left\{x\mid 2k\pi-\frac{\pi}{12}\leqslant x\leqslant 2k\pi+\frac{7\pi}{12}, k\in \Z\right\}\)}
    {\(\left\{x\mid 2k\pi+\frac{\pi}{6}\leqslant x\leqslant 2k\pi+\frac{5\pi}{6}, k\in \Z\right\}\)}
    {\(\left\{x\mid 2k\pi\leqslant x\leqslant 2k\pi+\frac{\pi}{6}\text{\ 或\ }2k\pi+\frac{5\pi}{6}\leqslant x\leqslant (2k+1)\pi, k\in \Z\right\}\)}

\end{problem}

\begin{answer}
A
\end{answer}

\begin{solution}
令 \(u=x-\frac{\pi}{4}\). 由 \(\sin u\geqslant\frac{1}{2}\) 得 \(2k\pi+\frac{\pi}{6}\leqslant u\leqslant2k\pi+\frac{5\pi}{6}\)，其中 \(k\in\Z\). 因而 \(2k\pi+\frac{5\pi}{12}\leqslant x\leqslant2k\pi+\frac{13\pi}{12}\)，故选 A.
\end{solution}

\begin{problem}
点 \((4,0)\) 关于直线 \(5x+4y+21=0\) 的对称点是（\quad）

\choices
    {\((-6,8)\)}
    {\((-8,-6)\)}
    {\((6,8)\)}
    {\((-6,-8)\)}

\end{problem}

\begin{answer}
D
\end{answer}

\begin{solution}
设点 \(P=(4,0)\)，直线为 \(5x+4y+21=0\). 令 \(D=5\cdot4+4\cdot0+21=41\)，且 \(5^2+4^2=41\). 点 \(P\) 关于该直线的对称点 \(P'\) 的坐标为 \(\left(4-\frac{2\cdot5D}{5^2+4^2},0-\frac{2\cdot4D}{5^2+4^2}\right)=(-6,-8)\)，故选 D.
\end{solution}

\begin{problem}
极坐标方程 \(4\sin 2\theta=3\) 表示的曲线是（\quad）

\choices
    {二条射线}
    {二条相交直线}
    {圆}
    {抛物线}

\end{problem}

\begin{answer}
B
\end{answer}

\begin{solution}
在极坐标中，\(\sin2\theta=\frac{2xy}{x^2+y^2}\)，所以原方程等价于 \(8xy=3(x^2+y^2)\)，即 \(3x^2-8xy+3y^2=0\). 这是过原点的二次齐次方程，其对应的斜率满足 \(3m^2-8m+3=0\)，有两个不同的实根 \(m=\frac{4\pm\sqrt{7}}{3}\)，故表示二条相交直线，选 B.
\end{solution}

\begin{problem}
由数字 \(0\)，\(1\)，\(2\)，\(3\)，\(4\)，\(5\) 组成无重复数字的六位数，其中个位数字小于十位数字的共有（\quad）

\choices
    {\(210\) 个}
    {\(300\) 个}
    {\(464\) 个}
    {\(600\) 个}

\end{problem}

\begin{answer}
B
\end{answer}

\begin{solution}
先选个位数字 \(u\) 和十位数字 \(v\)，要求 \(u<v\). 从 \(0\)，\(1\)，\(2\)，\(3\)，\(4\)，\(5\) 中选这样的有序数对共有 \(\mathrm C_6^2=15\) 个. 若 \(u=0\)，有 \(5\) 个数对，此时其余四个数字均非零，填入前四位有 \(4!=24\) 种；若 \(u\neq0\)，有 \(10\) 个数对，剩余四个数字中含有数字 \(0\)，填入前四位且首位不为零有 \(4!-3!=18\) 种. 因而共有 \(5\cdot24+10\cdot18=300\) 个，故选 B.
\end{solution}

\begin{problem}
如图是周期为 \(2\pi\) 的三角函数 \(y=f(x)\) 的图象，那么 \(f(x)\) 可以写成（\quad）

\bitmapfigure[max width=0.42\linewidth]{img/1991/sannan/q14_fig1.png}

\choices
    {\(\sin(1+x)\)}
    {\(\sin(-1-x)\)}
    {\(\sin(x-1)\)}
    {\(\sin(1-x)\)}

\end{problem}

\begin{answer}
D
\end{answer}

\begin{solution}
由图可见，\(f(0)>0\)，且曲线在 \(x=0\) 附近下降并在 \(x=1\) 处过 \(x\) 轴. 函数 \(y=\sin(1-x)\) 满足 \(\sin(1-0)=\sin1>0\)，在 \(x=0\) 处的变化方向为下降，并且 \(f(1)=0\)，其周期为 \(2\pi\)，与图象相符，故选 D.
\end{solution}

\begin{problem}
设命题甲为 \(\lg x^2=0\)；命题乙为 \(x=1\). 那么（\quad）

\choices
    {甲是乙的充分条件，但不是乙的必要条件}
    {甲是乙的必要条件，但不是乙的充分条件}
    {甲是乙的充要条件}
    {甲不是乙的充分条件，也不是乙的必要条件}

\end{problem}

\begin{answer}
B
\end{answer}

\begin{solution}
由 \(\lg x^2=0\) 得 \(x^2=1\)，所以命题甲等价于 \(x=1\) 或 \(x=-1\). 命题乙 \(x=1\) 能推出命题甲，但命题甲不能推出命题乙，因此甲是乙的必要条件，但不是乙的充分条件，故选 B.
\end{solution}

\begin{problem}
\(\left(\sqrt{x}-\frac{2}{\sqrt{x}}\right)^6\) 的展开式中常数项是（\quad）

\choices
    {\(-160\)}
    {\(-20\)}
    {\(20\)}
    {\(160\)}

\end{problem}

\begin{answer}
A
\end{answer}

\begin{solution}
展开式通项中，若从第二项 \(-\frac{2}{\sqrt{x}}\) 取 \(k\) 次，则该项的幂次为 \(x^{(6-k)/2}x^{-k/2}=x^{3-k}\). 常数项要求 \(k=3\)，故常数项为 \(\mathrm C_6^3(\sqrt{x})^3\left(-\frac{2}{\sqrt{x}}\right)^3=20\cdot(-8)=-160\)，选 A.
\end{solution}

\begin{problem}
体积相等的正方体，球，等边圆柱（即底面直径与母线相等的圆柱）的全面积分别为 \(S_1\)，\(S_2\)，\(S_3\)，那么它们的大小关系为（\quad）

\choices
    {\(S_1<S_2<S_3\)}
    {\(S_1<S_3<S_2\)}
    {\(S_2<S_3<S_1\)}
    {\(S_2<S_1<S_3\)}

\end{problem}

\begin{answer}
C
\end{answer}

\begin{solution}
设三者体积均为 \(V\). 正方体的全面积为 \(S_1=6V^{2/3}\)，故 \(S_1^3=216V^2\). 球的半径满足 \(\frac{4}{3}\pi R^3=V\)，所以 \(S_2=4\pi R^2\)，从而 \(S_2^3=36\pi V^2\). 等边圆柱设底面半径为 \(r\)，则高为 \(2r\)，由 \(2\pi r^3=V\)，其全面积为 \(S_3=6\pi r^2\)，从而 \(S_3^3=54\pi V^2\). 由于 \(36\pi<54\pi<216\)，所以 \(S_2<S_3<S_1\)，选 C.
\end{solution}

\begin{problem}
曲线 \(2y^2+3x+3=0\) 与曲线 \(x^2+y^2-4x-5=0\) 的公共点的个数是（\quad）

\choices
    {\(4\)}
    {\(3\)}
    {\(2\)}
    {\(1\)}

\end{problem}

\begin{answer}
D
\end{answer}

\begin{solution}
由第一条曲线得 \(y^2=\frac{-3x-3}{2}\)，所以必须有 \(x\leqslant-1\). 代入第二条曲线得 \(x^2+\frac{-3x-3}{2}-4x-5=0\)，即 \(2x^2-11x-13=0\)，解得 \(x=\frac{13}{2}\) 或 \(x=-1\). 由 \(x\leqslant-1\) 只能取 \(x=-1\)，代回得 \(y=0\)，故公共点只有一个，选 D.
\end{solution}

\section{填空题}

\begin{problem}
椭圆 \(9x^2+16y^2=144\) 的离心率为\fillinblank{}.
\end{problem}

\begin{answer}
\(\frac{\sqrt{7}}{4}\)
\end{answer}

\begin{solution}
将椭圆方程化为 \(\frac{x^2}{16}+\frac{y^2}{9}=1\)，所以长半轴 \(a=4\)，短半轴 \(b=3\)，焦半距 \(c=\sqrt{a^2-b^2}=\sqrt{7}\). 因而离心率为 \(e=\frac{c}{a}=\frac{\sqrt{7}}{4}\).
\end{solution}

\begin{problem}
设复数 \(z_1=2-\i\)，\(z_2=1-3\i\)，则复数 \(\frac{\i}{z_1}+\frac{\bar{z}_2}{5}\) 的虚部等于\fillinblank{}.
\end{problem}

\begin{answer}
\(1\)
\end{answer}

\begin{solution}
因为 \(\frac{\i}{z_1}=\frac{\i}{2-\i}=\frac{\i(2+\i)}{5}=\frac{-1+2\i}{5}\)，且 \(\frac{\bar z_2}{5}=\frac{1+3\i}{5}\)，所以 \(\frac{\i}{z_1}+\frac{\bar z_2}{5}=\i\)，其虚部为 \(1\).
\end{solution}

\begin{problem}
已知圆台的上、下底面半径分别为 \(r\)，\(2r\)，侧面积等于上、下底面积之和，则圆台的高为\fillinblank{}.
\end{problem}

\begin{answer}
\(\frac{4r}{3}\)
\end{answer}

\begin{solution}
设圆台的母线长为 \(l\)，高为 \(h\). 侧面积为 \(\pi(r+2r)l=3\pi rl\)，两底面面积之和为 \(\pi r^2+\pi(2r)^2=5\pi r^2\). 由题意得 \(l=\frac{5r}{3}\). 圆台的轴截面给出 \(l^2=h^2+(2r-r)^2=h^2+r^2\)，所以 \(h^2=\frac{25r^2}{9}-r^2=\frac{16r^2}{9}\). 因为 \(h>0\)，故 \(h=\frac{4r}{3}\).
\end{solution}

\begin{problem}
\(\displaystyle\lim_{n\to\infty}\frac{4n\cdot 2^n+1}{n\cdot 3^n-1}=\)\fillinblank{}.
\end{problem}

\begin{answer}
\(0\)
\end{answer}

\begin{solution}
分子、分母同除以 \(n\cdot3^n\)，得 \(\frac{4\left(\frac{2}{3}\right)^n+\frac{1}{n3^n}}{1-\frac{1}{n3^n}}\). 当 \(n\to\infty\) 时，\(\left(\frac{2}{3}\right)^n\to0\)，\(\frac{1}{n3^n}\to0\)，所以原极限为 \(0\).
\end{solution}

\begin{problem}
在体积为 \(V\) 的斜三棱柱 \(ABC-A'B'C'\) 中，已知 \(S\) 是侧棱 \(CC'\) 上的一点，过点 \(S\)，\(A\)，\(B\) 的截面截得的三棱锥的体积为 \(V_1\)，那么过点 \(S\)，\(A'\)，\(B'\) 的截面截得的三棱锥的体积为\fillinblank{}.
\end{problem}

\begin{answer}
\(\frac{V}{3}-V_1\)
\end{answer}

\begin{solution}
设两底面所在平面之间的距离为 \(H\)，并设 \(S\) 到平面 \(ABC\)，平面 \(A'B'C'\) 的距离分别为 \(h_1\)，\(h_2\). 因为两平面平行，且 \(S\) 在线段 \(CC'\) 上，所以 \(h_1+h_2=H\). 设底面三角形 \(ABC\) 的面积为 \(B\)，则棱柱体积为 \(V=BH\). 两个三棱锥的底面面积都为 \(B\)，故 \(V_1=\frac{1}{3}Bh_1\)，另一三棱锥体积为 \(V_2=\frac{1}{3}Bh_2\). 因而 \(V_1+V_2=\frac{1}{3}B(h_1+h_2)=\frac{V}{3}\)，所以 \(V_2=\frac{V}{3}-V_1\).
\end{solution}

\begin{problem}
设函数 \(f(x)=x^2+x+\frac{1}{2}\) 的定义域是 \(\{n,n+1\}\)（\(n\) 是自然数），那么在 \(f(x)\) 的值域中共有\fillinblank{} 个整数.
\end{problem}

\begin{answer}
\(0\)
\end{answer}

\begin{solution}
因为 \(n\) 是自然数，所以 \(n\) 与 \(n+1\) 都是整数. 于是 \(f(n)=n^2+n+\frac{1}{2}\)，\(f(n+1)=n^2+3n+\frac{5}{2}\)，二者都不是整数. 因此值域中整数的个数为 \(0\).
\end{solution}

\section{解答题}

\begin{problem}
已知 \(\alpha\)，\(\beta\) 为锐角，\(\cos \alpha=\frac{4}{5}\)，\(\tan(\alpha-\beta)=\frac{1}{3}\)，求 \(\cos \beta\) 的值.
\end{problem}

\begin{solution}
因为 \(\alpha\) 为锐角且 \(\cos\alpha=\frac{4}{5}\)，所以 \(\sin\alpha=\frac{3}{5}\)，从而 \(\tan\alpha=\frac{3}{4}\). 又因为 \(\alpha-\beta\in\left(-\frac{\pi}{2},\frac{\pi}{2}\right)\) 且 \(\tan(\alpha-\beta)=\frac{1}{3}>0\)，所以 \(\alpha>\beta\). 设 \(\tan\beta=t>0\)，则
\(\frac{\frac{3}{4}-t}{1+\frac{3}{4}t}=\frac{1}{3}\)，即 \(9-12t=4+3t\)，所以 \(t=\frac{1}{3}\). 因为 \(\beta\) 为锐角，\(\cos\beta=\frac{1}{\sqrt{1+\tan^2\beta}}=\frac{1}{\sqrt{1+\frac{1}{9}}}=\frac{3\sqrt{10}}{10}\).
\end{solution}

\begin{problem}
解不等式 \(\sqrt{5-4x-x^2}\geqslant x\).
\end{problem}

\begin{solution}
首先由根式有意义得 \(5-4x-x^2\geqslant0\)，即 \((x+2)^2\leqslant9\)，所以 \(-5\leqslant x\leqslant1\). 当 \(-5\leqslant x\leqslant0\) 时，左边非负而右边不大于零，不等式恒成立. 当 \(x>0\) 时，两边均非负，可以平方，得 \(5-4x-x^2\geqslant x^2\)，即 \(2x^2+4x-5\leqslant0\). 该不等式的两个根为 \(-1-\frac{\sqrt{14}}{2}\)，\(-1+\frac{\sqrt{14}}{2}\)，所以在 \(x>0\) 中的解为 \(0<x\leqslant-1+\frac{\sqrt{14}}{2}\). 合并得 \(x\in\left[-5,-1+\frac{\sqrt{14}}{2}\right]\).
\end{solution}

\begin{problem}
如图，已知直棱柱 \(ABC-A_1B_1C_1\) 中，\(\angle ACB=90^{\circ}\)，\(\angle BAC=30^{\circ}\)，\(BC=1\)，\(AA_1=\sqrt{6}\)，\(M\) 是 \(CC_1\) 的中点. 求证：\(AB_1\perp A_1M\).

\bitmapfigure[max width=0.30\linewidth]{img/1991/sannan/q27_fig1.png}
\end{problem}

\begin{solution}
在底面直角三角形 \(ABC\) 中，\(\angle A=30^{\circ}\)，\(\angle C=90^{\circ}\)，且 \(BC=1\)，所以 \(AB=2\)，\(AC=\sqrt{3}\). 建立空间直角坐标系，取 \(C=(0,0,0)\)，\(B=(1,0,0)\)，\(A=(0,\sqrt{3},0)\)，并令侧棱方向为 \(z\) 轴正方向. 因为棱柱为直棱柱且 \(AA_1=\sqrt{6}\)，有 \(B_1=(1,0,\sqrt{6})\)，\(A_1=(0,\sqrt{3},\sqrt{6})\)，\(C_1=(0,0,\sqrt{6})\)，\(M=\left(0,0,\frac{\sqrt{6}}{2}\right)\). 因而
\(\overrightarrow{AB_1}=(1,-\sqrt{3},\sqrt{6})\)，\(\overrightarrow{A_1M}=\left(0,-\sqrt{3},-\frac{\sqrt{6}}{2}\right)\). 于是 \(\overrightarrow{AB_1}\cdot\overrightarrow{A_1M}=0+3-3=0\)，所以 \(AB_1\perp A_1M\).
\end{solution}

\begin{problem}
设 \(\{a_n\}\) 是等差数列，\(a_1=1\)，\(S_n\) 是它的前 \(n\) 项和；\(\{b_n\}\) 是等比数列，其公比的绝对值小于 \(1\)，\(T_n\) 是它的前 \(n\) 项和. 如果 \(a_3=b_2\)，\(S_5=2T_2-6\)，\(\displaystyle\lim_{n\to\infty} T_n=9\)，求 \(\{a_n\}\)，\(\{b_n\}\) 的通项公式.
\end{problem}

\begin{solution}
设等差数列的公差为 \(d\)，等比数列的首项和公比分别为 \(b_1\)，\(q\). 则 \(a_3=1+2d\)，\(S_5=5(1+2d)=5a_3=5b_2=5b_1q\). 又由 \(\left|q\right|<1\) 及无穷等比数列前 \(n\) 项和的极限，\(\frac{b_1}{1-q}=9\)，即 \(b_1=9(1-q)\). 由 \(S_5=2T_2-6=2b_1(1+q)-6\)，得 \(5b_1q=2b_1(1+q)-6\)，即 \(b_1(3q-2)=-6\). 代入 \(b_1=9(1-q)\) 得 \(9(1-q)(3q-2)=-6\)，整理为 \(9q^2-15q+4=0\)，所以 \(q=\frac{4}{3}\) 或 \(q=\frac{1}{3}\). 由 \(\left|q\right|<1\) 取 \(q=\frac{1}{3}\)，从而 \(b_1=6\). 此时 \(b_2=2\)，故 \(a_3=2\)，即 \(1+2d=2\)，得 \(d=\frac{1}{2}\). 因而 \(a_n=1+\frac{1}{2}(n-1)\)，\(b_n=6\left(\frac{1}{3}\right)^{n-1}\).
\end{solution}

\begin{problem}
已知双曲线 \(C\) 的实半轴长与虚半轴的乘积为 \(\sqrt{3}\)，\(C\) 的两个焦点分别为 \(F_1\)，\(F_2\)，直线 \(l\) 过 \(F_2\) 且与直线 \(F_1F_2\) 的夹角满足 \(\tan \varphi=\frac{\sqrt{21}}{2}\)，\(l\) 与线段 \(F_1F_2\) 的垂直平分线的交点是 \(P\)，线段 \(PF_2\) 与双曲线 \(C\) 的交点为 \(Q\)，且 \(|PQ|:|QF_2|=2:1\). 求双曲线 \(C\) 的方程.
\end{problem}

\begin{solution}
设双曲线的标准方程为 \(\frac{x^2}{a^2}-\frac{y^2}{b^2}=1\)，焦点为 \(F_1=(-c,0)\)，\(F_2=(c,0)\)，其中 \(c^2=a^2+b^2\). 由直线 \(l\) 过 \(F_2\)，且与 \(x\) 轴的夹角满足 \(\tan\varphi=\frac{\sqrt{21}}{2}\)，它与垂直平分线 \(x=0\) 的交点可取为 \(P=\left(0,\frac{c\sqrt{21}}{2}\right)\). 因为 \(|PQ|:|QF_2|=2:1\)，点 \(Q\) 将线段 \(PF_2\) 按 \(2:1\) 分割，所以 \(Q=\left(\frac{2c}{3},\frac{c\sqrt{21}}{6}\right)\). 代入双曲线方程，并令 \(A=a^2\)，\(B=b^2\)，得
\(\frac{4c^2}{9A}-\frac{21c^2}{36B}=1\)，即 \((A+B)(16B-21A)=36AB\). 又由实半轴长与虚半轴长的乘积为 \(\sqrt{3}\)，有 \(AB=a^2b^2=3\). 展开上式得 \(16B^2-41AB-21A^2=0\)，令 \(r=\frac{B}{A}>0\)，则 \(16r^2-41r-21=0\)，所以 \(r=3\) 或负数，故 \(B=3A\). 结合 \(AB=3\) 得 \(A=1\)，\(B=3\)，所以双曲线方程为 \(x^2-\frac{y^2}{3}=1\).
\end{solution}

\begin{problem}
已知函数 \(f(x)=\frac{2^x-1}{2^x+1}\).
\begin{enumerate}
\item 证明：\(f(x)\) 在 \((-\infty,+\infty)\) 上是增函数；
\item 证明：对于任意不小于 \(3\) 的自然数 \(n\)，都有 \(f(n)>\frac{n}{n+1}\).
\end{enumerate}

\begin{solution}
\begin{enumerate}
\item 设 \(x_2>x_1\)，则 \(2^{x_2}>2^{x_1}>0\). 因而
\(f(x_2)-f(x_1)=\frac{2^{x_2}-1}{2^{x_2}+1}-\frac{2^{x_1}-1}{2^{x_1}+1}=\frac{2\left(2^{x_2}-2^{x_1}\right)}{\left(2^{x_2}+1\right)\left(2^{x_1}+1\right)}>0\). 所以 \(f(x)\) 在 \((-\infty,+\infty)\) 上是增函数.
\item 因为 \(2^n+1>0\) 且 \(n+1>0\)，不等式 \(f(n)>\frac{n}{n+1}\) 等价于 \((n+1)(2^n-1)>n(2^n+1)\)，即 \(2^n>2n+1\). 当 \(n=3\) 时，\(2^3=8>7=2\cdot3+1\). 若对某个 \(n\geqslant3\) 有 \(2^n>2n+1\)，则 \(2^{n+1}>4n+2>2n+3=2(n+1)+1\). 由数学归纳法，任意不小于 \(3\) 的自然数 \(n\) 都满足 \(2^n>2n+1\)，从而 \(f(n)>\frac{n}{n+1}\).
\end{enumerate}
\end{solution}
\end{problem}
